\section{Introduction}
\label{sec:introduction}

Zero-knowledge virtual machines (zkVMs) compile general-purpose
computation into succinct cryptographic proofs.  The Ziren
zkVM~\cite{ziren} is a MIPS-based zkVM built on the
Plonky3~\cite{plonky3} STARK toolkit, currently shipping with FRI as
its polynomial commitment scheme (PCS).  Although FRI offers
excellent verification time, modern multilinear PCS such as
WHIR~\cite{whir2024} promise simultaneously smaller proofs,
faster verification, and a more natural multilinear-extension (MLE)
view of trace data --- all of which are increasingly relevant as
zkVMs adopt hypercube-style architectures~\cite{sp1-hypercube}.

This paper documents the cryptographic and engineering work required
to migrate Ziren's STARK pipeline from FRI to a production-sound
WHIR-based pipeline.  The transition is non-trivial: the existing
``WHIR fast path'' shipped in Ziren is a structural skip of
permutation traces and quotient evaluation, with several unbound
soundness gaps.  We close all of them.

\subsection{Contributions}

\begin{enumerate}
  \item \textbf{Sumcheck-based zerocheck for transition constraints
        (\S\ref{sec:zerocheck-phase2a}).}  Replaces the quotient
        polynomial argument with a sumcheck protocol
        proving $\sum_{b \in \{0,1\}^m} \mathsf{eq}(r, b) \cdot
        C(b) = 0$, with provable soundness $O(d \cdot m / |\mathbb{F}|)$
        for constraint degree $d$.

  \item \textbf{Sumcheck-based LogUp-GKR for lookup
        arguments (\S\ref{sec:logup-gkr-fix}).}  We identify and
        fix a critical soundness bug in the initial folding-only
        implementation: the combine identity
        $(N(z\Vert 0)\cdot D(z\Vert 1) + D(z\Vert 0) \cdot
         N(z\Vert 1), D(z\Vert 0) \cdot D(z\Vert 1))$
        is degree-$2$ in each coordinate, which only matches the
        verifier's degree-$1$ multilinear $N_{k+1}, D_{k+1}$ at
        Boolean points.  Our replacement protocol uses a per-layer
        sumcheck reduction over a degree-$3$ round polynomial to
        ensure soundness at random points.

  \item \textbf{Late-binding WHIR adapter
        (\S\ref{sec:whir-late-binding}).}  The upstream WHIR
        \texttt{MultilinearPcs} adapter requires evaluation points
        \emph{at commit time}, which is incompatible with our
        protocol where the per-chip row coordinates
        $r_{\mathrm{row}}$ are derived \emph{after} commit via GKR.
        We construct an in-tree adapter built directly on
        \texttt{p3-whir} internals (\texttt{CommitmentWriter},
        \texttt{Verifier}, \texttt{EqStatement}) that exposes
        single-column and multi-column commit + post-commit
        \texttt{add\_late\_claim} + open API surfaces.  The
        cryptographic substrate is unchanged from upstream WHIR;
        only the transcript ordering differs.

  \item \textbf{Cross-binding (C.1) Brakedown spot-check
        protocol.}  In a dual-commit setting where FRI and WHIR
        both commit to the same trace, a malicious prover can
        commit divergent traces $A \neq B$.  We close this with a
        Brakedown-style spot-check argument~\cite{brakedown}: $K$
        random trace positions are sampled post-commit; the prover
        provides merkle paths against \emph{both} commits at each
        position.  Tampering on either side is detected with
        soundness $1 - (1 - 1/N)^K$.  For $K = 80$ and $N = 2^{17}$
        this gives $\sim 2^{-80}$ missed-detection probability.

  \item \textbf{Jagged late-binding construction.}  We combine
        Jagged~\cite{jagged} variable-height trace packing with the
        late-binding adapter via a sumcheck reduction that ties
        per-chip per-column row-MLE claims to a single dense
        commitment.  This collapses $\sim 570$ per-chip-per-column
        WHIR commits to a single shard-level commit, reducing
        proof size by a factor of $57$--$85\times$ in our
        evaluation (\S\ref{sec:performance}).

  \item \textbf{Recursion-circuit WHIR verifier scaffolding.}  We
        port the host-side verification logic into the
        recursion-circuit through a series of focused sub-systems:
        a \texttt{SumcheckVerify} chip (full DSL$\to$AIR pipeline,
        seven layers), a merkle-path verifier built atop the
        existing Poseidon2 chip, OOD evaluation primitives, and
        STIR query position derivation.  Each sub-system is
        independently testable and serves as a stepping stone
        toward in-circuit \texttt{verify\_whir\_pcs}.
\end{enumerate}

\subsection{Implementation and Open-Source Release}

All host-side prover and verifier code is implemented in the open-
source Ziren codebase, opt-in via two environment variables:
\texttt{ZIREN\_USE\_WHIR=1} (selects the WHIR fast path with
zerocheck and LogUp-GKR), and \texttt{ZIREN\_LATE\_BINDING=jagged}
(selects the jagged late-binding path).  A third flag
\texttt{ZIREN\_CROSS\_BINDING=1} activates the cross-binding spot-
check protocol.  Defaults preserve the original FRI pipeline so
existing deployments are unaffected.

We test each cryptographic mechanism via end-to-end round-trip and
tamper-rejection unit tests, plus end-to-end smoke tests on the
fibonacci, JSON, and keccak workloads.  At the time of writing,
$24$ late-binding-related tests, $9$ zerocheck and LogUp-GKR tests,
$18$ recursion-circuit WHIR verifier tests, and the existing $34$
recursion-core tests all pass.  Compressed-proof verification of
each smoke-tested workload completes successfully.

\subsection{Paper Outline}

\S\ref{sec:preliminaries} reviews the necessary preliminaries
(STARK, sumcheck, WHIR, LogUp-GKR, Plonky3 architecture).
\S\ref{sec:background} sets the engineering context, summarises
the Plonky3 upgrade work, and tabulates implementation status.
\S\ref{sec:plonky3-upgrade} details the recursion-circuit FRI
verifier upgrade required by upstream Plonky3 protocol changes.
\S\ref{sec:soundness-fix} formalises the threat model of the
existing WHIR fast path.
\S\S\ref{sec:zerocheck-phase2a}--\ref{sec:whir-late-binding}
present the three new constructions: zerocheck (Phase 2a),
LogUp-GKR (Phase 2b), and late-binding WHIR with cross-binding
and jagged packing (Phase 2c+).
\S\ref{sec:performance} reports performance measurements.
\S\ref{sec:related-work} positions the work against prior
constructions.  \S\ref{sec:conclusion} concludes with future
directions.  Appendices document detailed implementation status
and a session-by-session changelog.
